\chapter{Ingresos energéticos y empresas públicas}

\titular{La transferencia del Fondo Mexicano del Petróleo cae \textbf{20,7\,\% real} y las
ventas de Pemex crecen. La relación fiscal entre el Estado y su petrolera se desplaza: el
Estado recibe menos renta y \textbf{aporta más capital}. El Ramo 18 pasa de 138\,307,4 a
267\,439,1 millones de pesos, y \textbf{263\,476,3 de esos son una sola partida: compra de
títulos y valores}.}

\quecambio{El Fondo Mexicano del Petróleo transfiere 232\,630,4 millones contra 279\,766,8 en
la iniciativa anterior. Las ventas de bienes y servicios de Pemex se estiman en 971\,677,2
millones y las de la Comisión Federal de Electricidad en 535\,477,2. El precio supuesto del
barril baja de 62,0 a 54,9 dólares y la plataforma de exportación de 616 a 521 mil barriles
diarios, una caída de 15,4\,\%.}

\porqueimporta{\textbf{La caída de la renta petrolera no es un accidente de precio: es
sobre todo de volumen exportado.} Y la contrapartida está del lado del gasto, en una inyección
de capital que casi duplica el Ramo 18. El Estado está sustituyendo ingreso petrolero por
inversión financiera en la empresa que lo genera.}

\section{Declaración de perímetro y de colocación}

\marginal{qué está y qué no}
Este capítulo trata la \emph{relación fiscal} entre el Estado y las empresas públicas, que es
un objeto de la Parte I: la renta que el Estado recibe, las metas de balance que les fija y el
capital que les aporta. \textbf{El presupuesto de gasto de Pemex y de la Comisión Federal de
Electricidad no está aquí}: vive en el capítulo 4, dentro de los agregados, y en el 8, en la
inversión. Se dice para que el lector no lo busque en estas páginas.

\section{La renta petrolera}

El artículo 1o.\ recoge la renta petrolera en dos lugares. El numeral 9.97, transferencias del
Fondo Mexicano del Petróleo para la Estabilización y el Desarrollo, vale \textbf{232\,630,4
millones de pesos}: 47\,136,4 menos en nominal y 60\,565,2 menos en términos reales que en la
iniciativa anterior, una caída de \textbf{20,7\,\%}. Y el numeral 1.18.01, impuesto por la
actividad de exploración y extracción de hidrocarburos, vale 7\,070,4 millones.

Sumados a los ingresos propios de Pemex por venta de bienes y servicios ---971\,677,2
millones--- el bloque energético del artículo 1o.\ se distribuye así: el Estado recibe
directamente \textbf{239\,700,8 millones} y la empresa retiene \textbf{971\,677,2}.

\section{Una ausencia que hay que declarar}

\textbf{La iniciativa de 2026 no contiene un artículo que fije o sustituya la tasa del derecho
petrolero.} Entre 2020 y 2023 el artículo 22 de la Ley de Ingresos sustituía año con año la
tasa del derecho por la utilidad compartida ---65, 58, 54 y 40 por ciento en aquellos
ejercicios--- y ese artículo era el lugar donde se leía la política de renta petrolera del
año.

Se buscó y no está. \textbf{Este documento no supone una tasa}: declara que la iniciativa no
la contiene y que, por tanto, la política de renta petrolera de 2026 no se puede leer de la
Ley de Ingresos como se leía hasta 2023. Es una pérdida de trazabilidad, y es exactamente el
tipo de cosa que un lector del paquete debería poder verificar en un artículo.

\section{Las empresas: metas y capital}

El proyecto de decreto fija a Pemex y a la Comisión Federal de Electricidad sus metas de
balance financiero, y la iniciativa de ingresos les autoriza techos de endeudamiento propios:
\textbf{160\,619,6 millones internos y 5\,342,1 millones de dólares externos} para Pemex;
\textbf{8\,764,2 millones internos y 969,0 millones de dólares} para la Comisión.

\marginal{el capital, no el subsidio}
Y luego está el Ramo 18. Pasa de 138\,307,4 a \textbf{267\,439,1 millones de pesos}, un
aumento de 84,5\,\% nominal. \textbf{El 98,5\,\% del ramo es una sola partida específica, la
73903, «adquisición de otros valores»}: 263\,476,3 millones en 2026 contra 136\,210,3 en 2025.
Es capítulo 7000, inversión financiera. No es subsidio de operación ni transferencia
corriente: \textbf{es compra de capital}.

La cifra tiene su reflejo en el marco de mediano plazo: la inversión financiera del sector
público pasa de 0,4 a \textbf{0,7\,\% del PIB} entre 2025 y 2026, y vuelve a 0,3 en 2027. Es
un movimiento de un solo año.

\textbf{Esta aportación y el gasto propio de la empresa no se suman.} El neteo del Anexo 1
---1\,553\,113,1 millones--- descuenta subsidios, transferencias y apoyos fiscales a las
entidades de control directo y a las empresas del Estado, precisamente para que el gasto neto
total no cuente dos veces el mismo peso. Cualquier lectura que sume el Ramo 18 y el
presupuesto de Pemex está contando de más.

\section{Implicaciones}

La política energética del paquete 2026 se lee mejor en el gasto que en el ingreso, y eso ya
es un hallazgo. El Estado proyecta recibir menos renta ---20,7\,\% real menos del Fondo--- y
comprometer más capital, con una inyección que duplica el Ramo 18 y que el propio escenario
oficial no repite en 2027.

Queda una pregunta que este documento no puede responder con el paquete en la mano:
\textbf{qué parte de los 263\,476,3 millones corresponde a amortización de pasivos de la
empresa y qué parte a inversión productiva.} La partida no lo distingue y la exposición de
motivos del proyecto de presupuesto, que sería el lugar natural para explicarlo, lleva cinco
ejercicios caída en el servidor de la Secretaría.

\begin{ficha}{ingresos energéticos y empresas públicas}
\fila{Perímetro}{Renta petrolera del artículo 1o.\ (numerales 9.97 y 1.18.01) más ingresos
propios de las dos empresas (numeral 7.72). El gasto de las empresas queda fuera y se trata en
los capítulos 4 y 8.}
\fila{Línea de comparación}{\textbf{Línea P:} iniciativa 2026 contra iniciativa 2025.}
\fila{Deflactor}{Del PIB, 4,8\,\%, CGPE 2026.}
\fila{Año de los pesos}{Corrientes; las diferencias reales en pesos de 2026.}
\fila{PIB y añada}{38\,715,9 miles de millones, CGPE 2026.}
\fila{Denominador demográfico}{No aplica.}
\fila{Fuentes}{ILIF 2026 e ILIF 2025, artículo 1o.\ y artículos 4o.\ y 6o.; CGPE 2026 Anexos
II.5 y III.2; proyecto de decreto, Anexo 1; base del proyecto por partida específica.}
\fila{Qué no puede afirmarse}{La tasa del derecho petrolero de 2026: la iniciativa no la fija
y no se supone. La composición de la aportación de capital entre amortización de pasivos e
inversión: la partida no lo distingue. Y el balance financiero proyectado de cada empresa
contra su meta, que exigiría los estados financieros y no está en el paquete.}
\end{ficha}
\clearpage
